% Packages, verbatim style, and title block.
% Loaded by main.tex; do not put \documentclass or \begin{document} here.

% NeurIPS 2024 style: US letter, 5.5in x 9in text block, 10pt Times, title between
% rules. [preprint] prints the authors (the default mode anonymises them and adds
% line numbers); [nonatbib] keeps the numeric IEEEtran references.
\usepackage[preprint,nonatbib]{neurips_2024}

% Language and fonts. T5 must be the last (default) encoding so that Vietnamese
% IRIs and labels (Nguyễn_Công_Phượng, Thể_loại:…) typeset with pdflatex.
% neurips_2024 asks for Times (ptm) and Helvetica (phv), which have no T5 glyphs;
% TeX Gyre Termes and Heros are the same designs with T5 support. Latin Modern
% stays as the typewriter font for IRIs and code.
\usepackage[utf8]{inputenc}
\usepackage[T1,T5]{fontenc}
\usepackage{lmodern}
\usepackage{tgtermes}
\usepackage{tgheros}
\usepackage[vietnamese,english]{babel}

% Useful packages
\usepackage{amsmath, amssymb}
\usepackage{graphicx}
\usepackage{adjustbox}
\usepackage{xcolor}
\usepackage{float}
\usepackage[font=small,labelfont=bf]{caption}
\usepackage{subcaption}
\usepackage{microtype}

\usepackage{booktabs}
\usepackage{multirow}
\usepackage{tabularx}
\usepackage{array}

% Verbatim blocks (SPARQL, Turtle, curl). listings cannot handle UTF-8
% Vietnamese characters, fancyvrb passes them through to the T5 font.
% neurips_2024 makes \footnotesize 9pt; 8pt keeps the longest lines (~95
% characters) inside the 5.5in text block. samepage keeps a block on one page.
\usepackage{fancyvrb}
\fvset{fontsize=\fontsize{8pt}{9.5pt}\selectfont, samepage=true, frame=single,
       framesep=1.5mm, rulecolor=\color{gray!60}}
% Floating listings keep a verbatim block on one page: \begin{listing}…\caption{…}\end{listing}
\newfloat{listing}{t}{lol}
\floatname{listing}{Listing}

% "approximately" tilde. \sim is a mathrel, so bare $\sim$5 gets thick
% relation spacing on both sides; \mathord strips it.
\newcommand{\apx}{\ensuremath{\mathord{\sim}}}

% Vocabulary terms and IRIs in running text: \term{vio:playedFor}.
% Typewriter fonts do not hyphenate by default; long IRIs would overrun the
% margin, so allow hyphenation in \term and give paragraphs extra stretch.
\newcommand{\term}[1]{\text{\ttfamily\hyphenchar\font=45\relax #1}}
\setlength{\emergencystretch}{3em}

% Placeholder the group must fill in before submission (prints in red).
\newcommand{\fillin}[1]{\textcolor{red}{[#1]}}

% --- TikZ for schematic figures ---
\usepackage{tikz}
\usetikzlibrary{arrows.meta, calc, positioning, backgrounds, shapes.geometric}
% Accent colours: project vocabulary (blue), inferred triples (orange),
% external datasets and asserted edges (gray).
\definecolor{accblue}{HTML}{2A78D6}
\definecolor{accorange}{HTML}{EB6834}
\definecolor{accgray}{HTML}{8A8984}
% RDFS graph figure (figures/rdfs_graph.tex): instance nodes, property nodes and
% domain/range edges, entailed triples
\colorlet{rdfsdraw}{accblue}
\colorlet{rdfsfill}{accblue!6}
\colorlet{rdfsprop}{accblue!75!black}
\colorlet{rdfsaxiom}{accorange}

\usepackage[colorlinks=true, allcolors=blue]{hyperref}

\title{Building a Vietnamese DBpedia:\\ Linked Open Data from Vietnamese Wikipedia}

% NeurIPS author block: \And separates authors on one line; the first line of
% each block is bold, the next lines are the student ID.
\author{%
  Ha Thi Kim Phung \\ 20242460M
  \And
  Do Ngoc Trung \\ 20232311M
  \And
  Bui Duc Phan Hoang \\ 20252004M
  \And
  Do Hoang Dung \\ 20241184M
}

% Footnote at the bottom of page 1, where NeurIPS prints the conference notice.
\makeatletter
\renewcommand{\@noticestring}{%
  Semantic Web course project, Topic~2: \emph{Build a DBpedia version for Vietnamese language}.
  Group~23, lecturer Dr.~Do Ba Lam. Hanoi University of Science and Technology, October 2026.%
}
\makeatother
